% ============================================================================
% Fichier  : partie4/P04_C15_forensique_reseau.tex
% Titre    : Forensique Réseau Opérationnelle
% Auteur   : MINKA MI NGUIDJOÏ Thierry Emmanuel
% Version  : 1.0 -- Juin 2026
% Statut   : RÉDIGÉ
% Sources  : Ch. 9 (Wireshark -- bases), Ch. 14 (redondance forensique),
%            M5_Investigation_Numerique_PNC.docx (BLOC 4 réseau),
%            affaires WOURI (BEC), BAOBAB (MoMo transfrontalier)
% Positionnement : Ch. 9 a posé les bases Wireshark (filtres de base, cas
%                  WOURI). Ce chapitre approfondit : capture positionnée,
%                  protocoles chiffrés, exfiltration DNS, analyse de logs
%                  pare-feu/proxy, corrélation multi-sources.
% ============================================================================

\chapter{Forensique Réseau Opérationnelle}
\label{chap:for-res}

\epigraph{%
    « Le réseau ne ment jamais. Il ne se souvient pas de tout,
    mais ce dont il se souvient est vrai. »%
}{--- Principe de la forensique réseau}

\niveauChapitre{Expert}{%
    Maîtrise du Chapitre~\ref{chap:acq} (Wireshark bases)
    et du Chapitre~\ref{chap:for-sys} (timeline, redondance)
    indispensable. Protocoles TCP/IP, DNS, HTTP/S requis.%
}

\begin{boxobjectifs}
\begin{itemize}
    \item Positionner une sonde de capture dans différentes topologies
          réseau et comprendre ce que chaque position capture (et rate).
    \item Analyser en profondeur une capture \texttt{.pcap}~:
          flux TCP, indicateurs d'exfiltration, tunnels DNS,
          beaconing de malware.
    \item Reconstruire des sessions applicatives depuis une capture~:
          emails SMTP, transferts HTTP, commandes FTP.
    \item Analyser des logs de pare-feu, proxy et DNS pour reconstituer
          l'activité réseau en l'absence de capture complète.
    \item Identifier les marqueurs de communications C2
          (\textit{Command and Control}) dans le trafic chiffré.
    \item Appliquer la corrélation temporelle multi-sources~:
          capture réseau + logs système + timeline (Ch.~\ref{chap:for-sys}).
\end{itemize}
\end{boxobjectifs}

\bigskip

% ============================================================================
\section*{Positionnement dans le cours}
% ============================================================================

Le Chapitre~\ref{chap:acq} a introduit Wireshark avec ses filtres
de base et le cas WOURI. Ce chapitre va plus loin~: comment capturer
intelligemment dans un réseau réel, comment analyser le trafic chiffré,
comment détecter l'exfiltration de données, et comment reconstituer
l'activité réseau quand on n'a pas de capture (mais des logs).

\separateur

% ============================================================================
\section{Positionnement de la Sonde~: ce qu'on Voit et ce qu'on Rate}
\label{sec:positionnement-sonde}
% ============================================================================

\subsection{Topologies et points de capture}

La valeur d'une capture réseau dépend \emph{entièrement} de la position
de la sonde dans le réseau. Capturer au mauvais endroit, c'est ne rien voir.

\begin{tableaucours}{p{3.5cm}p{4.0cm}p{4.5cm}}{%
    \tableauHeader{Position} &
    \tableauHeader{Ce qu'on capture} &
    \tableauHeader{Ce qu'on rate}
}
\textbf{Port SPAN / Mirror} sur le switch principal &
    Tout le trafic qui traverse le switch (nord-sud ET est-ouest) &
    Trafic interne entre machines sur le même segment si non mirrored \\
\textbf{TAP réseau} sur le lien WAN &
    Tout le trafic entrant/sortant de l'entreprise vers Internet &
    Trafic interne (LAN to LAN)~; communications intra-entreprise \\
\textbf{Agent hôte} (sur la machine suspecte) &
    Tout le trafic de \emph{cette} machine, même chiffré avant
    le réseau &
    Trafic des autres machines \\
\textbf{Pare-feu / proxy logs} (sans capture) &
    Connexions autorisées/refusées (IP, port, octets) &
    Contenu des connexions~; trafic qui ne passe pas par le
    proxy (UDP, ICMP) \\
\end{tableaucours}
\vspace{-0.8em}
\begin{center}{\small\itshape Tableau~15.1 --- Positions de capture et couverture}\end{center}

\subsection{Capture sur un switch managé~: SPAN port}

\begin{boxoutils}{Configuration d'un SPAN port (Cisco IOS)}
\begin{lstlisting}[language=bash_forensique]
! Configuration d'une session SPAN sur Cisco IOS
! Port monitore : GigabitEthernet0/1 (machine suspecte)
! Port destination : GigabitEthernet0/24 (poste forensique)

Switch(config)# monitor session 1 source interface Gi0/1 both
Switch(config)# monitor session 1 destination interface Gi0/24

! Verifier la configuration
Switch# show monitor session 1
! Resultat :
! Session 1
! Type           : Local Session
! Source Ports   : Both : Gi0/1
! Destination Ports : Gi0/24

! Sur le poste forensique (connecte sur Gi0/24) :
tcpdump -i eth0 -w /output/capture_suspect.pcap
! Ou avec ring buffer (fichiers de 100 Mo max) :
tcpdump -i eth0 -C 100 -W 50 -w /output/capture_%Y%m%d_%H%M%S.pcap
\end{lstlisting}
\end{boxoutils}

\begin{boiteRemarque}{Wireshark live vs tcpdump + analyse offline}
En contexte forensique, la bonne pratique est de capturer avec
\texttt{tcpdump} (léger, consomme peu de ressources) en production,
puis d'analyser avec Wireshark offline sur le \texttt{.pcap} résultant.
Lancer Wireshark avec son interface graphique directement sur une
machine en production peut introduire une charge réseau supplémentaire
et des artefacts dans la capture elle-même.
\end{boiteRemarque}

% ============================================================================
\section{Analyse Approfondie des Flux TCP}
\label{sec:analyse-flux}
% ============================================================================

\subsection{Suivre un flux complet~: Follow TCP Stream}

La fonction \textit{Follow TCP Stream} de Wireshark reconstitue la
session complète d'un échange TCP en réassemblant les segments dans
l'ordre. C'est l'outil le plus puissant pour comprendre rapidement
ce qu'une connexion a échangé.

\begin{boxoutils}{Wireshark -- Follow TCP Stream et export}
\begin{lstlisting}[language=bash_forensique]
# En ligne de commande avec tshark (equivalent Wireshark CLI)
# Extraire tous les flux TCP et les afficher
tshark -r capture.pcap -q -z "follow,tcp,ascii,0"

# Lister les conversations TCP avec volumes
tshark -r capture.pcap -q -z "conv,tcp"
# Resultat :
#                    |      <-      |      ->      |    Total
# 10.0.0.15:49212 <-> 185.220.XX.XX:4444
#                         1476       14832          16308 bytes

# Extraire les objets HTTP (fichiers telecharges)
tshark -r capture.pcap --export-objects "http,/output/http_objects/"

# Extraire les objets SMB (fichiers partages reseau)
tshark -r capture.pcap --export-objects "smb,/output/smb_objects/"
# Utile pour detecter exfiltration via partages reseau

# Calculer le hash SHA-256 de chaque objet extrait
find /output/http_objects/ -type f -exec sha256sum {} \; \
    > /output/http_objects_hashes.txt
# Puis soumettre sur VirusTotal (API ou interface web)
\end{lstlisting}
\end{boxoutils}

\subsection{Métriques d'anomalie réseau}

\begin{tableaucours}{p{3.5cm}p{4.0cm}p{4.5cm}}{%
    \tableauHeader{Métrique} &
    \tableauHeader{Valeur normale} &
    \tableauHeader{Anomalie (indicateur suspect)}
}
Durée d'une session HTTP & 0.1--2~secondes & $>$~60~secondes~:
    session persistante anormale \\
Volume de réponse HTTP & Variable & Requête~\texttt{GET /index.html}
    de 200~octets reçoit 50~Mo~: exfiltration \\
Ratio upload/download & $\ll$~1 (navigation) & Ratio proche de~1~:
    exfiltration~; ratio $\gg$~1~: upload de données \\
Intervalles de connexion & Irréguliers & Intervalles fixes (30~s,
    60~s, 300~s)~: beaconing de malware \\
Entropie du contenu & Faible (texte/HTML) & Haute (binaire, chiffré)~:
    contenu encodé ou chiffré \\
\end{tableaucours}
\vspace{-0.8em}
\begin{center}{\small\itshape Tableau~15.2 --- Métriques d'anomalie réseau}\end{center}

% ============================================================================
\section{Détection d'Exfiltration de Données}
\label{sec:exfiltration}
% ============================================================================

\subsection{Exfiltration HTTP/HTTPS~: le vecteur le plus courant}

\begin{boxoutils}{Détecter une exfiltration HTTP depuis une capture}
\begin{lstlisting}[language=bash_forensique]
# Requetes POST avec corps volumineux (upload de donnees)
tshark -r capture.pcap -Y "http.request.method == POST" \
    -T fields -e frame.time -e ip.src -e ip.dst \
    -e http.host -e http.content_length

# Resultat suspect :
# 2026-03-13 02:19:05  10.0.0.15  185.220.XX.XX
#   update.windowscdn.net  content_length: 45231400
# -> 45 Mo envoyes vers un domaine qui semble legitime
#    mais est enregistre il y a 2 jours (CDN de facade)

# Calculer le volume par destination (top exfiltration targets)
tshark -r capture.pcap -q -z "ip_hosts,tree"
# Alternatif : Statistics --> IP Addresses dans Wireshark GUI

# Detecter contenu base64 dans les parametres URL
tshark -r capture.pcap -Y "http" \
    -T fields -e http.request.uri | \
    grep -oE "[A-Za-z0-9+/]{30,}={0,2}"
# Base64 de 30+ caracteres dans les URLs = encodage de donnees voles
\end{lstlisting}
\end{boxoutils}

\subsection{Tunneling DNS~: exfiltration via le port~53}

Le \termecle{tunneling DNS}\index{tunneling DNS} exploite le protocole
DNS pour exfiltrer des données ou établir un canal C2 (Command \& Control).
Son avantage~: le port~53 UDP est ouvert sur presque tous les pare-feu.

\begin{boxoutils}{Détecter le tunneling DNS}
\begin{lstlisting}[language=bash_forensique]
# Statistiques des requetes DNS par domaine (top queried)
tshark -r capture.pcap -Y "dns.qry.type == A" \
    -T fields -e dns.qry.name | sort | uniq -c | sort -rn | head -20

# Indicateurs de tunneling DNS :
# 1. Sous-domaines tres longs (encodage base32/64 dans le label)
tshark -r capture.pcap -Y "dns" -T fields -e dns.qry.name | \
    awk 'length($0) > 50'
# Resultat suspect :
# aGVsbG8gd29ybGQgdGhpcyBpcyBhIHRlc3Q.malware.domain.com
# -> sous-domaine de 60 caracteres = payload encode

# 2. Volume de requetes DNS anormalement eleve pour un meme domaine
tshark -r capture.pcap -Y "dns.qry.name contains malware.domain.com" \
    -T fields -e frame.time -e dns.qry.name | wc -l
# 847 requetes en 10 minutes vers le meme domaine = tunnel

# 3. Reponses DNS TXT avec contenu inhabituel (canal C2 entrant)
tshark -r capture.pcap -Y "dns.resp.type == 16" \
    -T fields -e dns.txt
# Si les TXT ressemblent a du binaire encode : C2 via DNS TXT

# Outil specialise : dnscat2-detector ou dnstap-analysis
# Pour la forensique offline : analyser les logs du resolver DNS
\end{lstlisting}
\end{boxoutils}

\begin{boxCRO}
\textbf{Analyse CRO --- Tunneling DNS et détection}

$\mathcal{C}$ (Confidentialité pour l'attaquant)~: le tunnel DNS chiffre
généralement son contenu~; même si le tunnel est détecté, le contenu
exfiltré peut être inaccessible sans la clé de déchiffrement.

$\mathcal{R}$ (Fiabilité de la détection)~: les indicateurs
statistiques (longueur de requête, volume, TXT records) sont des
\emph{indicateurs} de tunneling, pas des preuves directes. Un domaine
CDN légitime peut générer des sous-domaines longs. Documenter la
convergence de plusieurs indicateurs.

$\mathcal{O}$ (Opposabilité)~: la capture pcap avec les requêtes DNS
anormales constitue une preuve directe de l'activité réseau, admissible
selon les mêmes principes que tout autre artefact réseau.

\cro{$\downarrow$ contenu inaccessible}{$\uparrow$ convergence indicateurs}{$\uparrow\uparrow$ capture pcap directe}
\end{boxCRO}

\subsection{Beaconing~: détecter un implant actif}

Le \termecle{beaconing}\index{beaconing} est le comportement caractéristique
d'un malware qui contacte périodiquement son serveur C2 pour recevoir
des instructions. La signature~: des connexions à intervalles
réguliers vers la même destination.

\begin{lstlisting}[language=bash_forensique,
    caption={Détection de beaconing dans une capture}]
# Extraire les timestamps de connexions vers une IP suspecte
tshark -r capture.pcap \
    -Y "ip.dst == 185.220.XX.XX && tcp.flags.syn == 1" \
    -T fields -e frame.time_epoch | sort -n > /tmp/beacon_times.txt

# Calculer les intervalles entre connexions (en secondes)
awk 'NR>1{printf "%.0f\n", $1-prev} {prev=$1}' \
    /tmp/beacon_times.txt | sort | uniq -c | sort -rn

# Resultat type (beaconing sur 300 secondes) :
#   47 300    <- 47 intervalles de exactement 300 secondes
#    2 299
#    1 301
# -> malware avec beacon toutes les 5 minutes, jitter +/- 1s
# Comparaison : trafic normal -> distribution aleatoire des intervalles
\end{lstlisting}

% ============================================================================
\section{Reconstruction de Sessions Applicatives}
\label{sec:reconstruction-sessions}
% ============================================================================

\subsection{Reconstituer un email SMTP}

\begin{boxoutils}{Extraction et reconstruction d'email SMTP}
\begin{lstlisting}[language=bash_forensique]
# Extraire le contenu SMTP brut
tshark -r capture.pcap -Y "smtp" \
    -T fields -e smtp.req.command -e smtp.req.parameter

# Follow TCP Stream sur une session SMTP (non chiffree)
tshark -r capture.pcap -q -z "follow,tcp,ascii,5"
# Resultat type :
# EHLO attacker.ru
# MAIL FROM:<contact@medequip-portugal.com>
# RCPT TO:<finance@sotradi.cm>
# DATA
# From: Service Comptabilite <contact@medequip-portugal.com>
# Subject: Facture mensuelle equipements
# [corps du message avec piece jointe frauduleuse]

# Extraire les pieces jointes SMTP (format MIME)
# Wireshark GUI : File --> Export Objects --> IMF (Internet Mail)
# CLI :
tshark -r capture.pcap --export-objects "imf,/output/smtp_objects/"
# Puis analyser la piece jointe :
file /output/smtp_objects/attachment.pdf
sha256sum /output/smtp_objects/attachment.pdf
# Verifier sur VirusTotal
\end{lstlisting}
\end{boxoutils}

\begin{boxscenario}{Cas WOURI --- Reconstruction complète de l'email BEC}
La capture réseau de SOTRADI~SA (\texttt{sotradi\_capture.pcap})
contient la session SMTP du phishing. Analyse~:

\textbf{Étape~1 --- Identifier la session SMTP suspecte~:}
\begin{lstlisting}[language=bash_forensique]
# Filtrer les sessions SMTP entrantes (serveur mail SOTRADI sur TCP 25)
tshark -r sotradi_capture.pcap -Y "smtp && tcp.dstport == 25"
# Session identifiee : frame 14872-14950, IP source 91.XXX.XX.XX
\end{lstlisting}

\textbf{Étape~2 --- Analyser les en-têtes SMTP~:}
\begin{lstlisting}[language=bash_forensique]
tshark -r sotradi_capture.pcap -q -z "follow,tcp,ascii,14872"
# En-tetes reveles :
# MAIL FROM: contact@medequip-portugal.com
# X-Originating-IP: 91.XXX.XX.XX    <- IP reelle (serveur russe)
# Received: from mail.yandex.ru      <- origine reelle
# Date: 2026-01-15 14:28:03
\end{lstlisting}

\textbf{Constatation forensique~:}
L'email de la session TCP~5 (frame~14872) affiche l'adresse expéditrice
\texttt{contact@medequip-portugal.com} dans le champ \texttt{MAIL FROM}
mais révèle dans le champ \texttt{X-Originating-IP} l'IP
\texttt{91.XXX.XX.XX}, enregistrée en Russie selon les bases WHOIS
consultées le 20~janvier~2026. L'en-tête \texttt{Received} confirme
le transit par les serveurs Yandex.ru.
\end{boxscenario}

\subsection{Analyse de trafic FTP~: transferts de fichiers en clair}

\begin{lstlisting}[language=bash_forensique,
    caption={Reconstruction de transferts FTP depuis une capture}]
# FTP en clair (port 21 controle, 20 donnees) : tout visible
tshark -r capture.pcap -Y "ftp" \
    -T fields -e frame.time -e ftp.request.command \
    -e ftp.request.arg

# Resultat type :
# 02:20:15  USER   admin
# 02:20:16  PASS   admin123  <- mot de passe en clair !
# 02:20:17  CWD    /confidential/RH/
# 02:20:18  RETR   salaires_2026.xlsx  <- telechargement fichier confidentiel

# Extraire les fichiers transferes via FTP-DATA (port 20)
tshark -r capture.pcap --export-objects "ftp-data,/output/ftp_objects/"

# Note : FTPS (FTP over TLS) chiffre le contenu mais les commandes
# (USER, PASS, fichiers) restent potentiellement visibles si TLS
# est mal configure (certificat auto-signe, version obsolete)
\end{lstlisting}

% ============================================================================
\section{Analyse des Logs sans Capture Complète}
\label{sec:analyse-logs}
% ============================================================================

En l'absence de capture pcap (la situation la plus courante en
investigation~: le réseau n'était pas monitoré au moment de l'incident),
les logs de périphériques réseau deviennent la seule source d'information
réseau. Ils sont moins riches, mais souvent disponibles.

\subsection{Logs de pare-feu}

\begin{boxoutils}{Analyse de logs de pare-feu (format syslog/Palo Alto)}
\begin{lstlisting}[language=bash_forensique]
# Logs de pare-feu au format syslog standard
# Structure typique : timestamp src_ip dst_ip src_port dst_port action

# Extraire toutes les connexions vers une IP suspecte
grep "185.220.XX.XX" /var/log/firewall.log | \
    grep "ACCEPT\|ALLOW" | \
    awk '{print $1,$2,$3,$4,$5,$6}' | sort -k3 -n

# Calculer le volume par connexion (bytes transferes)
grep "185.220.XX.XX" /var/log/firewall.log | \
    awk '{bytes+=$NF} END {print "Total bytes:", bytes}'

# Connexions refusees (tentatives d'intrusion ou scan)
grep "DENY\|BLOCK\|DROP" /var/log/firewall.log | \
    awk '{print $5}' | sort | uniq -c | sort -rn | head -20
# Top IP sources bloquees : peut identifier un scanner

# Detecter le beaconing dans les logs pare-feu
grep "185.220.XX.XX" /var/log/firewall.log | \
    awk '{print $1}' | \
    awk 'BEGIN{FS=":"}{printf "%s%s%s\n",$1,$2,$3}' | \
    sort | uniq -c | sort -rn
# Connexions recurrentes a la meme minute -> beaconing probable
\end{lstlisting}
\end{boxoutils}

\subsection{Logs de proxy web}

\begin{boxoutils}{Analyse de logs de proxy Squid/Bluecoat}
\begin{lstlisting}[language=bash_forensique]
# Format Squid access.log :
# timestamp elapsed client action/code bytes method url

# Sites les plus consultés (volume)
awk '{print $7}' /var/log/squid/access.log | \
    awk -F/ '{print $3}' | sort | uniq -c | sort -rn | head -20

# Telechargements volumineux (potentielle exfiltration)
awk '$5 > 10000000 {print $3,$7,$5}' /var/log/squid/access.log
# -> connections > 10 Mo vers des URLs externes

# User-Agents inhabituels (malware, outils d'attaque)
awk '{print $NF}' /var/log/squid/access.log | sort | uniq -c | \
    sort -rn | grep -iv "Mozilla\|Chrome\|Firefox\|Safari"
# Resultat suspect : "python-requests/2.28.0" -> script automatise
#                    "curl/7.74.0" -> ligne de commande
#                    "Go-http-client/1.1" -> outil Go

# Connexions hors horaires de bureau
awk '$1 >= 1678668000 && $1 <= 1678675200' /var/log/squid/access.log
# (timestamps UNIX pour 00h00-02h00 le 13 mars 2026)
\end{lstlisting}
\end{boxoutils}

\subsection{Logs DNS~: reconstituer l'activité réseau par les noms}

\begin{lstlisting}[language=bash_forensique,
    caption={Analyse de logs DNS (BIND / Windows DNS)}]
# Logs de requetes DNS (BIND format)
grep "QUERY" /var/log/named/query.log | \
    awk '{print $5}' | sort | uniq -c | sort -rn | head -30

# Domaines enregistres recemment (DGA / domaines d'attaque)
# Les DGA (Domain Generation Algorithms) produisent des noms aleatoires
grep "QUERY" /var/log/named/query.log | \
    awk '{print $5}' | \
    awk -F. '{
        n = split($0, a, ".");
        label = a[n-1] "." a[n];  # domaine principal
        sub_label = a[1];          # premier sous-domaine
        if (length(sub_label) > 15) print $0  # sous-domaine long
    }'

# Volume de requetes par client (qui fait beaucoup de DNS ?)
grep "QUERY" /var/log/named/query.log | \
    awk '{print $4}' | sort | uniq -c | sort -rn | head -10
# Si une machine fait 10x plus de requetes que les autres : suspect

# Synchroniser DNS logs avec timeline systeme (Ch. 14)
# Timestamp DNS 02:14:07 -> resolution medequip-portugal.com
# Timestamp MFT 02:17:44 -> execution enc_all.ps1
# Correlation : la resolution DNS precede l'execution de 3 min 37 sec
\end{lstlisting}

% ============================================================================
\section{Corrélation Temporelle Multi-Sources}
\label{sec:correlation-temporelle}
% ============================================================================

La puissance de la forensique réseau réside dans la \textbf{corrélation}
entre les sources réseau et les sources système. Chaque source seule
est un fragment~; la corrélation révèle l'image complète.

\begin{tableaucours}{p{2.5cm}p{4.0cm}p{5.0cm}}{%
    \tableauHeader{Source} &
    \tableauHeader{Information réseau} &
    \tableauHeader{Corrélation avec Ch.~\ref{chap:for-sys}}
}
\textbf{Capture pcap} & IP, ports, contenu des sessions &
    Corrélation temporelle~: connexion TCP à 02h14:09 $\leftrightarrow$
    entrée journald SSH Accepted à 02h14:09 \\
\textbf{Logs pare-feu} & Flux autorisés/refusés, volumes &
    Corrélation volumétrique~: 45~Mo sortants à 02h19
    $\leftrightarrow$ \texttt{\$UsnJrnl} accès 847~fichiers à 02h18 \\
\textbf{Logs DNS} & Noms résolus, clients, timestamps &
    Corrélation causale~: résolution \texttt{malware.domain.com}
    à 02h13~$\to$ connexion vers cette IP à 02h14 \\
\textbf{Logs proxy} & URLs, User-Agents, volumes &
    Corrélation comportementale~: User-Agent \texttt{python-requests}
    depuis machine du comptable = inhabituel \\
\end{tableaucours}
\vspace{-0.8em}
\begin{center}{\small\itshape Tableau~15.3 --- Corrélation temporelle multi-sources}\end{center}

\begin{boxscenario}{Timeline réseau + système~: Opération MVOM complète}
En fusionnant la timeline système (Ch.~\ref{chap:for-sys}) avec les
sources réseau de ce chapitre, la séquence complète de l'attaque émerge~:

\begin{lstlisting}[language=bash_forensique]
# Timeline fusionnee (source en brackets)
02:13:45  [DNS log]     Resolution: update.windowscdn.net -> 185.220.XX.XX
02:14:02  [journald]    SSH Failed password from 185.220.XX.XX (x7)
02:14:09  [journald]    SSH Accepted password root from 185.220.XX.XX
02:14:09  [pcap]        TCP SYN dst=185.220.XX.XX:4444 (C2 channel)
02:16:50  [$UsnJrnl]    FileCreate: enc_all.ps1
02:17:40  [UserAssist]  powershell.exe executed
02:17:44  [Prefetch]    enc_all.ps1 executed
02:18:31  [$UsnJrnl]    847 fichiers .docx DataOverwrite (chiffrement)
02:19:05  [firewall]    ALLOW 10.0.0.15 -> 185.220.XX.XX:443 45.2MB
02:19:05  [proxy log]   POST https://update.windowscdn.net 45231400 bytes
02:19:48  [pcap]        TCP FIN 185.220.XX.XX:4444 (deconnexion C2)
\end{lstlisting}

\textbf{Reconstruction narrative}~: L'attaquant a d'abord effectué
une résolution DNS pour préparer sa connexion (02h13). Il a obtenu
un accès root par brute-force SSH (02h14). Il a immédiatement
ouvert un canal C2 (port~4444, 02h14) par lequel il a téléchargé
et exécuté le ransomware (02h16--02h18). Il a exfiltré 45~Mo de
données (vraisemblablement les fichiers avant chiffrement, 02h19)
avant de fermer la connexion C2 (02h19:48).
\end{boxscenario}

\begin{boxjuris}
\textbf{Admissibilité des logs réseau en droit camerounais}

Les logs de pare-feu, proxy et serveurs DNS constituent des ``données
informatiques'' au sens de l'art.~4 de la \loicam{2010/012}.
Leur admissibilité en tant que preuves est soumise aux mêmes
conditions que tout artefact numérique~:
\begin{enumerate}
    \item Origine légale~: les logs doivent avoir été conservés dans
          le cadre normal de l'exploitation du réseau (art.~27
          \loicam{2024/017}~: obligation de sécurité et traçabilité)~;
    \item Intégrité~: hash SHA-256 calculé au moment de la saisie~;
    \item Chaîne de custody documentée depuis le serveur jusqu'au
          tribunal~;
    \item Expert judiciaire agréé pour attester de l'authenticité
          et de l'interprétation technique.
\end{enumerate}
\end{boxjuris}

\separateur

\begin{boxresume}
\textbf{Ce qu'il faut retenir de ce chapitre~:}
\begin{itemize}
    \item \textbf{Position de la sonde d'abord}~: port SPAN (tout
          le switch)~; TAP WAN (tout le trafic Internet)~;
          agent hôte (tout le trafic de la machine, chiffré compris).
          Capturer au mauvais endroit~= ne rien voir.

    \item \textbf{Follow TCP Stream}~: reconstruit la session complète.
          \texttt{--export-objects http} extrait les fichiers transférés
          avec leurs hash. Calculer le SHA-256 de chaque objet et
          soumettre sur VirusTotal.

    \item \textbf{Exfiltration~:} détecter par le volume (POST
          volumineux)~; l'entropie du contenu (contenu chiffré)~;
          le ratio upload/download anormal.

    \item \textbf{Tunneling DNS}~: sous-domaines de~$>$~50~caractères~;
          volume anormal de requêtes vers un même domaine~;
          réponses TXT avec contenu binaire.

    \item \textbf{Beaconing}~: intervalles réguliers de connexion vers
          la même destination. Calculer les intervalles et leur
          distribution~: un malware produit une distribution
          quasi-normale centrée sur son intervalle~; le trafic
          légitime est aléatoire.

    \item \textbf{Logs sans capture}~: logs pare-feu (flux, volumes)~;
          proxy (URLs, User-Agents)~; DNS (domaines résolus) constituent
          des sources complémentaires indispensables quand pas de pcap.

    \item \textbf{Corrélation temporelle}~: la valeur de la forensique
          réseau vient de la croisé avec la timeline système
          (Ch.~\ref{chap:for-sys}). La convergence temporelle entre
          log DNS~$\to$ log pare-feu~$\to$ journald~$\to$
          \texttt{\$UsnJrnl} est une preuve bien plus robuste que
          chaque source isolée.

    \item \textbf{Admissibilité}~: logs réseau~= données informatiques
          (art.~4 Loi~2010/012)~; mêmes exigences~: intégrité (hash),
          custody, expert agréé.
\end{itemize}
\end{boxresume}

\bigskip

\begin{boxexo}[title={\ding{45}\quad Exercice~15.1 --- Position de sonde \niveaubase}]
Une entreprise de Yaoundé suspecte une exfiltration de données depuis
un poste de travail du département comptabilité. Le réseau comprend~:
un routeur Internet, un switch principal, 40~postes de travail Windows,
un serveur de fichiers, et un serveur de messagerie.
\begin{enumerate}[(a)]
    \item Où positionneriez-vous la sonde de capture pour voir tout
          le trafic Internet du poste suspect~?
    \item Pourrait-on capturer le trafic interne entre le poste suspect
          et le serveur de fichiers avec cette même sonde~? Si non, comment~?
    \item Si vous n'avez pas accès au switch pour un SPAN port, quelles
          alternatives avez-vous pour capturer le trafic~?
    \item À quel moment de l'incident est-il trop tard pour une capture
          réseau (et quelles sources de substitution utiliser)~?
\end{enumerate}
\end{boxexo}

\begin{boxexo}[title={\ding{45}\quad Exercice~15.2 --- Analyse de capture \niveaumoyen}]
La capture \texttt{incident\_baobab.pcap} de l'Opération BAOBAB
(Chapitre~\ref{chap:norm-glob}) est disponible. \texttt{tshark -r
incident\_baobab.pcap -q -z "conv,tcp"} révèle~:
\begin{lstlisting}[language=bash_forensique]
# Extrait conversations TCP :
10.0.0.42:52341 <-> 196.200.XX.XX:443   847 pkts  2.3 MB  <->  4.1 KB
10.0.0.42:52411 <-> 196.200.XX.XX:443   848 pkts  2.3 MB  <->  4.2 KB
10.0.0.42:52478 <-> 196.200.XX.XX:443   847 pkts  2.3 MB  <->  4.1 KB
[...] (total : 47 conversations similaires)
\end{lstlisting}
\begin{enumerate}[(a)]
    \item Quelle anomalie identifiez-vous dans ces conversations~?
          À quoi ressemble ce pattern~?
    \item Calculez l'intervalle approximatif entre les connexions.
    \item Quelles commandes tshark utiliseriez-vous pour~: (i)~extraire
          les intervalles entre connexions~; (ii)~vérifier si le
          contenu est chiffré~; (iii)~identifier le User-Agent si
          ces connexions sont HTTP/S~?
    \item Rédigez la constatation forensique correspondante (sans
          interprétation) et l'analyse avec niveau de confiance.
\end{enumerate}
\end{boxexo}

\begin{boxexo}[title={\ding{45}\quad Exercice~15.3 --- Corrélation multi-sources \niveauexpert}]
Vous disposez des sources suivantes pour l'Opération WOURI~:
\begin{itemize}
    \item Capture \texttt{sotradi.pcap} (trafic réseau 14--15~janvier 2026)~;
    \item Logs pare-feu Palo Alto (30~jours)~;
    \item Logs proxy Bluecoat (30~jours)~;
    \item Logs DNS interne (30~jours)~;
    \item Image forensique E01 du poste du comptable~;
    \item Dump RAM du serveur de messagerie.
\end{itemize}
\begin{enumerate}[(a)]
    \item Construisez un plan d'investigation réseau~: dans quel ordre
          analyser ces sources~? Quelle question spécifique posez-vous
          à chaque source~?
    \item Comment corréleriez-vous temporellement les logs DNS avec
          le contenu SMTP de la capture pcap pour établir une
          chronologie de l'attaque BEC~?
    \item Appliquez la grille CRO à la décision d'extraire et d'analyser
          le contenu des pièces jointes SMTP identifiées dans la
          capture~: qui peut y accéder~? Quelle base légale cite-t-on~?
    \item Rédigez la section Constatations réseau de votre rapport,
          couvrant la période 14h00--15h00 du 15~janvier~2026.
\end{enumerate}
\end{boxexo}

\bigskip

\begin{boxinfo}
\textbf{Transition.} Le Chapitre~\ref{chap:anti-for} (Anti-Forensique
et Contre-Mesures) traite l'autre côté du miroir~: les techniques
qu'un attaquant sophistiqué utilise pour effacer les traces dans les
systèmes (Ch.~\ref{chap:for-sys}) et le réseau (ce chapitre)~---
et comment les détecter malgré tout. La connaissance des techniques
d'anti-forensique renforce la capacité à les contrer ET à les documenter
quand elles ont été utilisées.
\end{boxinfo}
